\par\noindent\textbf{(1)} 取全局相位不影响强度的波片矩阵J=diag(e\textasciicircum{}\{iδ/2\},e\textasciicircum{}\{-iδ/2\})，输入E0(cosα,sinα)变为Eout=E0(e\textasciicircum{}\{iδ/2\}cosα,e\textasciicircum{}\{-iδ/2\}sinα)。两个正交场幅不变而相位差改变量δ，故纯无损波片保持总强度。
\par\noindent\textbf{(2)} 采用Q=|Ex|²−|Ey|²、U=2Re(ExEy*)、V=2Im(ExEy*)，归一Stokes为(Q,U,V)=(cos2α,sin2αcosδ,sin2αsinδ)，满足Q²+U²+V²=1。α=π/4且δ=±π/2时Q=U=0、|V|=1，为两种相反手性的圆偏振。
\par\noindent\textbf{(3)} 理想x检偏器只透过Ex，故Iout/Iin=|e\textasciicircum{}\{iδ/2\}cosα|²=cos²α，与该固定轴检偏器不显含δ。这说明单用波片主轴上的一片检偏器不能测出相位迟滞，必须改变波片轴角或加入另一个分析器。
\par\noindent\textbf{(4)} 半波片δ=π令输入线偏振α在Stokes赤道上变为反射方向−α（相对于波片主轴），因此实验室中若波片主轴与x夹角β，输出偏振方位为2β−α。波片转β引起输出方位转2β，不能误写成同角度旋转。
